\section{Introduction}
%% Giorgio

%% Paul
% Hard Switched DC-DC Converters: What they are, why they were used, and what are their problems regarding switching losses and EMI.
% Softswitched DC-DC Converters: What they are, how they work, and their advantages in terms of reduced switching losses and EMI. But they come with increased complexity and therefore weight, which is crucial in satellite applications.
% Something like this (gemini): 
For the DC DC converter different topologies exist, which can be broadly categorized into hard-switched and soft-switched (resonant) converters. 
The primary motivation for investigating resonant converter topologies in ion propulsion is the relationship between switching frequency and converter system mass. In traditional power electronic circuits, using hard-switched devices, switching losses are directly proportional to the operating frequency. While increasing frequency is the standard method for reducing the size of magnetic components like inductors, and therefore their weight, it leads to power dissipation in hard-switched systems where voltage and current overlap during transitions.Theoretical research~\cite{bellar1998review}, suggests that by integrating a resonant network, the converter can achieve Zero Voltage Switching (ZVS) or Zero Current Switching (ZCS). This effectively shapes the switching waveforms so that transitions occur only when the voltage or current is zero, hence the power disspated during switching is minimized. The key advantage is hence, decoupling efficiency from frequency. This is especially relevant for converters for satellite applications, such as ion thrusters, where the core engineering challenge is maximizing the effective specific power ($kW/kg$) of components~\cite{gao2017soft}. On the other hand, soft-switched converters need additional circuitry to achieve ZVS or ZCS, which increases the overall system complexity and mass. Hence a comparison of the two topologies in terms of efficiency and mass is crucial for determining the optimal design for ion thruster PPUs.